\section{The normalization of the real Bessel period}
\label{sec:bessel-periods}

The unipotent estimates ensure convergence in the required weight range. We still need to compare the period formed with a definite compact type with the Hilbert norm of the ambient discrete series. Interlacing selects the model for which that comparison is made. The resulting isometry and decay
estimate enter the formal-degree calculation of the next
chapter. Definition~\ref{hyp:real-bessel-measures} specifies the marked negative vector, its positive orthogonal complement and the unitary character of the flag radical. The quotient norm in Proposition~\ref{prop:real-period-norm} uses those measures.

\begin{paramdefinition}[Real forms and period measures]\label{hyp:real-bessel-measures}
Let \(n=k+2r+1\), let \(G=U(k+r,r+1)\), and choose the flag
basis for \(G\). Write \(N\) for the unipotent radical of the parabolic stabilizing the isotropic flag. In the middle space of signature \((k,1)\), choose a negative vector \(z\) for the Hermitian form \(h\). Its compact stabilizer and the Bessel subgroup are specified by
\[
      h(z,z)=-T_z<0,\qquad W=z^\perp,\qquad S=U(W)N .
\]
In particular \(W\) is positive definite.

The chosen
unitary character \(\psi\) of \(N\) is nontrivial on each of the
simple coordinates of the Bessel datum. The last coordinate pairs
with \(z\). The compact group preserves \(\psi\). Its action on a fixed
Haar measure on \(N\) has a positive real modulus, a continuous
homomorphism from \(U(W)\) to \(\mathbb R_{>0}\). Its image is
compact and hence is trivial. We use \(dh\,du\) on \(S=U(W)N\)
with these choices of \(dh\) and \(du\).
\end{paramdefinition}

\subsection{Interlacing}

The compact stabilizer acts on the target of the period. The relative model therefore uses its dual.

Let \(\mu_1>\cdots>\mu_n\) and \(u_1>\cdots>u_k\) be the
Harish-Chandra parameters of the large-group discrete series and
compact module. Here \(k\) is even. The parameters of that dual in decreasing order are \(-u_k,\ldots,-u_1\). In the classical normalization of the
real Gan--Gross--Prasad theorem, the anisotropic line in the
complement has sign \((-1)^k=+1\). Negate both Hermitian forms
to reach this normalization. The matrix groups are unchanged
and both packets exchange their positive and negative labels.
A diagonal change of the isotropic basis normalizes the simple
character coefficients. Starting at the last coordinate, solve
in turn for each consecutive ratio of the diagonal scalars.

\begin{proposition}[Interlacing on the prescribed packet]
Use Definition~\ref{hyp:real-bessel-measures} with \(k\) even and
\(\mu_k,\ldots,\mu_n<u_k\). After negating the two Hermitian
forms, take the first \(k-1\) entries negative and the tail
alternating from a negative entry. For this signed packet the tempered Bessel criterion holds precisely when 
\[
 u_i>\mu_i>u_{i+1}\qquad(1\le i<k).
\]

\end{proposition}

\begin{proof}
For a discrete parameter, Atobe's labeling gives component value
\((-1)^{i-1}\) at a positive entry in position \(i\), and \((-1)^i\)
at a negative entry. The two ranks have opposite parity. In the
epsilon criterion each factor associated with \(\mu_i,u_j\) is
therefore the root number of a character of odd exponent
\(2(\mu_i-u_j)\). This root number is \(-1\) for a positive exponent
and \(1\) for a negative exponent, in the normalization of Atobe's
Proposition~2.3. The required component values by the tempered
Gan--Gross--Prasad theorem are
\[
 (-1)^{\#\{j:u_j<\mu_i\}},\qquad
 (-1)^{\#\{i:\mu_i>u_j\}}
\]
at the large-group position \(i\) and the compact-dual position
\(k+1-j\), respectively.

After the sign change the first \(k-1\) large entries are negative
and the tail alternates negative and positive. If
\(u_i>\mu_i>u_{i+1}\) for \(1\le i<k\) and the tail lies below
\(u_k\), the two counts are
\[
 \#\{j:u_j<\mu_i\}=\begin{cases}k-i&i<k,\\0&i\ge k,\end{cases}
 \qquad \#\{i:\mu_i>u_j\}=j-1.
\]
Since \(k\) is even, the large component values are \((-1)^i\)
before the tail and \(1\) on the tail. This matches the required labels, since the
tail starts at the even position \(k\). The negative definite dual
has component value \((-1)^{k+1-j}=(-1)^{j-1}\), and exactly \(j-1\)
large parameters lie above \(u_j\).

Conversely, let \(n_j=\#\{i:\mu_i>u_j\}\). The compact label
requires \(n_j\equiv j-1\pmod2\). With the tail fixed below
\(u_k\), at most \(k-1\) entries are available, and
\[
   n_{j+1}-n_j\ge1\quad(1\le j<k),\qquad n_k-n_1\le k-1.
\]
Every difference is therefore one. This proves the strict interlacing.
These calculations use \href{https://eprints.lib.hokudai.ac.jp/repo/huscap/all/84508/Adv.%20math%20363_106984.pdf}{Atobe, Theorem 2.1 and Proposition 2.3},
\href{https://doi.org/10.1215/00127094-2022-0018}{Xue's tempered Bessel theorem},
and the conventions of \href{https://arxiv.org/html/2610.29483v1}{Chen, Li and Ren, Theorem 1.1 and Section 2}.

\end{proof}

\begin{construction}
\begin{tabular}{@{}lll@{}}\toprule
Entry & Count in the epsilon criterion & Packet sign\\\midrule
\(\mu_i,\ i<k\) & \(\#\{j:u_j<\mu_i\}=k-i\) & \((-1)^i\)\\
\(\mu_i,\ i\ge k\) & \(\#\{j:u_j<\mu_i\}=0\) & \(1\)\\
\(u_j\) at dual position \(k+1-j\) & \(\#\{i:\mu_i>u_j\}=j-1\) & \((-1)^{j-1}\)\\\bottomrule
\end{tabular}
\caption{Component signs for even \(k\) under strict interlacing, with the remaining large-group entries below \(u_k\). The compact-dual counts determine each gap between adjacent \(u_j\).}\label{fig:interlacing-sign-counts}
\end{construction}

\subsection{From a rank-one period to the big-cell norm}

Interlacing identifies the compact type. Its occurrence alone gives no scale for the integral, so we compute the rank-one operator using the previously fixed Haar measures.

Write \(\sigma\) for the unitary discrete series and \(\pi\) for
the compact representation. The period operator is
\[
 B(x,y)=\int_{U(W)\times N}
 \overline{\psi(u)}\langle y,\sigma(hu)x\rangle\pi(h)^{-1}\,dh\,du.
\]
\begin{proposition}[Period factorization and quotient norm]\label{prop:real-period-norm}
Use Definition~\ref{hyp:real-bessel-measures} and opposite-sign gaps
large enough that \(h/(8\sqrt{2(r+1)})-1>2rn^2\), as in
Proposition~\ref{prop:bessel-sign-gap-decay}. Then

\[
 B(x,y)=L(x)L(y)^*
\]
 on the smooth slot, with \(L\) unique up to unit phase when the period is
nonzero. For a finite sum \(F\) of compact-finite coefficients with
definite-type operator factors, let \(s\) be the sum of their
contractions against the two Bessel maps. In the specified Haar
measures the squared big-cell norm is \(c_{\rm meas}|s|^2/d_\sigma\).
Here \(d_\sigma\) is the formal degree and \(c_{\rm meas}\) depends
only on the measures and the cell setting.
\end{proposition}

\begin{proof}
\proofstep{Positivity and the scale of the period.}

It is absolutely convergent for smooth \(x\) and compact-finite
\(y\) in the range of gaps used here. We describe the decay
argument below. Invariance of Haar gives equivariance in the
first variable. The form on \(\sigma\otimes\pi^\vee\) obtained
by pairing the operator against compact vectors is positive.
Indeed, for a tensor \(v\) of compact-finite vectors, set
\(\rho=\bar\psi(\sigma\otimes\pi^\vee)\) on \(S\). For right
F\o lner sets \(A\) in this unimodular amenable group,
\[
 |A|^{-1}\left\|\int_A\rho(a)v\,da\right\|^2
 =\int_S \frac{|A\cap As^{-1}|}{|A|}
                   \langle v,\rho(s)v\rangle\,ds.
\]
The overlap factor tends to one and is bounded by one, so
absolute integrability permits dominated convergence. The limit
is the diagonal of the period form and is nonnegative. This
argument applies to finite sums of compact-finite tensors.

By Bessel uniqueness the form factors through a space of
equivariant functionals of dimension at most one. For the nonzero form \(\mathcal B\), fix \(v_0\) with
\(\mathcal B(v_0,v_0)>0\) and put
\(b(v)=\mathcal B(v,v_0)/\sqrt{\mathcal B(v_0,v_0)}\),
using the convention linear in \(v\).
By uniqueness, each functional \(\mathcal B(\cdot,v')\)
is proportional to \(b\). Evaluating at \(v_0\) and using
Hermitian symmetry fixes its multiplier:
\[
                \mathcal B(v,v')=b(v)\overline{b(v')}.
\]
Under the tensor-dual correspondence \(b\) defines the
map \(L:\sigma\to\pi\). The rank-one tensor on the right
then gives \(B(x,y)=L(x)L(y)^*\) in the convention of
the period operator. Another choice of \(v_0\) changes the
functional by a unit phase.

For the zero form take \(L=0\). The map is
nonzero in the distinguished packet by
\href{https://arxiv.org/pdf/1506.01452#page=181}{Beuzart--Plessis, Theorem 7.2.1}.
The canonical integral there agrees with absolute integration
in the current range. In its smooth dilation average, undo the
Haar change of variables and use that the dilation test has
integral one. This is justified by absolute convergence and the
bounded range of the dilation variable. Density of compact-finite
vectors and continuity give the factorization on the smooth slot.

\proofstep{Contraction by Schur orthogonality.}

The remaining scale comes from integrating against an ambient matrix coefficient and hence from the formal degree.


Take two compact-finite coefficients
\[
 F(g)=\langle x,\sigma(g)y\rangle X,\qquad
 F'(g)=\langle x',\sigma(g)y'\rangle X'.
\]
Let \(P(p)\) pair the two-sided translate of \(F\) with \(F'\)
on the group. Write
\[
 p=(h_L,u_L,h_R,u_R),\qquad s_L=u_Lh_L,\qquad s_R=u_Rh_R.
\]
Schur orthogonality gives
\[
 P(p)=\frac{\langle y',\sigma(s_R)y\rangle
                 \langle\sigma(s_L)x,x'\rangle}{d_\sigma}
      \langle\pi(h_L)X\pi'(h_R)^{-1},X'\rangle_{\rm HS},
\]
where the Hilbert--Schmidt pairing is linear in its first slot.
The two scalar coefficients of the period variables are absolutely
integrable. Invert the left variables and conjugate the right
unipotent by its compact variable to use the rank-one periods.
If \(\xi(p)\) is the product character of the two sides, then
\[
 \int \xi(p)P(p)\,dp
    =d_\sigma^{-1}\langle L(x),XL'(y)\rangle
                \overline{\langle L(x'),X'L'(y')\rangle}.
\]
After the two measure-preserving changes of variables, the
operator inside the pairing with \(X'\) is
\[
       L(x')L(x)^*\,X\,L'(y)L'(y')^* .
\]
Pairing this rank-one operator against \(X'\) gives the two scalar
contractions in the formula. In a sum of such coefficients, write \(s\) for
the sum of their scalar contractions against \(L,L'\).
The pairing on the diagonal is \(|s|^2/d_\sigma\).

\proofstep{Regularization at positive frequencies.}

Disintegrate Haar measure on the group into the product of the period measures and the measure \(d\mu_{\rm cell}(q)\) on the open cell, where \(q\) denotes its coordinates. At two positive frequencies,
regularize the translated group pairing by smooth averages of
the two positive dilations \(a(t_i)\), with compactly supported
weights of integral one. The simple coordinates are multiplied
by \(t_i\). The Haar power \((t_1t_2)^\kappa\) in the averaging
formula cancels after the substitutions in the period variables.


For the Fourier phase at frequency \(\tau_i\) and simple coordinate \(\lambda_i\), use
\[
 (1-\partial_{t_i}^2)e^{it_i\tau_i\lambda_i}
 =(1+\tau_i^2\lambda_i^2)e^{it_i\tau_i\lambda_i}.
\]
Integration by parts on the compactly supported weight gives
a factor \((1+\tau_i^2\lambda_i^2)^{-1}\) on each side. The Xi bound gives
an integrable majorant after these four derivatives and makes
the regularized pairing continuous in the positive frequencies.
When the translated pairing is integrable without the weights,
Fubini and those substitutions give the ordinary pairing.

\proofstep{The pointwise norm identity.}

The averaged big-cell norm gives the same regularized pairing.
Expand the square under smooth, compactly supported product tests
in \((\tau_1,\tau_2)\). For each one-dimensional test \(\varphi\in C_c^\infty\),
in the Fourier convention with phase \(e^{it\lambda}\),
\[
 (1+\lambda^2)|\widehat\varphi(\lambda)|
 =|\widehat{\varphi-\varphi''}(\lambda)|
 \le\|\varphi-\varphi''\|_1 .
\]
Applying this in each coordinate gives the product of the two decay factors. Disintegrate one
of the two period integrals using the group Haar formula, and use
unimodularity to move the quotient representative. The absolute
coefficient majorant and the two Xi integrals justify the exchanges.
The tests identify the squared norm and the regularized pairing
\(I_\tau\) at almost every positive frequency pair. The single period
converges absolutely at almost every cell point by the same Haar
formula applied to the norm of a coefficient. Its value there is
continuous in the frequencies. Approaching \(\tau\) by pairs
\(\tau_n\) from the equality set gives by Fatou
\[
   \|J_\tau F\|_2^2\le\liminf_n\|J_{\tau_n}F\|_2^2=I_\tau .
\]

For a coefficient with five
derivatives, changing the dilations changes it continuously in the
four-derivative seminorm in that bound. Denote the coefficient dilation by \(F\mapsto F_t\)
and the weighted \(L^2\) period by \(J_\tau F\).
The covariance and the bound obtained from Fatou give
\[
 \|J_{t\tau}F\|_2=(t_1t_2)^{-\kappa/2}\|J_\tau F_t\|_2,
 \qquad
 \big|\|J_\tau F_t\|_2-\|J_\tau F\|_2\big|
        \le \|J_\tau(F_t-F)\|_2 .
\]
As \(t\to(1,1)\), the right side tends to zero by the
four-derivative seminorm bound. Both sides of the
almost-everywhere squared-norm identity are now continuous in
the positive frequency quadrant, giving equality at \((1,1)\). For arbitrary fixed
Haar measures the result is
\[
  \int \|JF(g_q)\|_{\rm HS}^2\,d\mu_{\rm cell}(q)
       = c_{\rm meas}|s|^2/d_\sigma,
\]
where \(c_{\rm meas}>0\) depends on the measures and the cell
setting. It is independent of the representation in that setting.
This is the scale used in Section~\ref{sec:local-amplitudes}.

\end{proof}

\subsection{Discrete-series translation and decay}

Period factorization uses a gap condition. Translation within the open chamber supplies the necessary decay across the weight range.

Fix the compact Cartan and its ordered Borel. The translation
input used here is the tensor identity for cohomological
induction with all shifts in the same open chamber. A Borel
weight-line filtration then induces in a single good-range degree
and gives the translated discrete series with their weight
multiplicities. This is the discrete series case of \href{https://www.mathunion.org/fileadmin/ICM/Proceedings/ICM1978.2/ICM1978.2.ocr.pdf}{Zuckerman's coherent translation, Theorem B}.
For the exterior standard powers and the dual standard the weights
have multiplicity one. The shifts here remain ordered with distinct
infinitesimal characters. The center splits the filtration. The
signed Hermitian tensor form is orthogonal across these summands by
Schur's lemma. Nondegeneracy gives a nonzero real multiple of each
unitary form, which can be normalized to a sign.

\begin{proposition}[Decay in terms of the opposite-sign gap]\label{prop:bessel-sign-gap-decay}
When all opposite-sign gaps are at least \(h\), the coefficient of vectors \(v,w\) with sufficiently many derivatives satisfies 
\[
 |\langle v,\sigma(g)w\rangle|\ll_{v,w,m}\|g\|_{\rm HS}^{-m}\qquad(m\in\mathbb Z,\ 0\le m<h).
\]

For a fixed \(m\), the seminorms of the two vectors can be taken of
order \(m(n-1)\). This bound also applies to a fixed derivative of
the coefficient after differentiating the vectors.
\end{proposition}

\begin{proof}
\proofstep{Regular shifts and coefficient estimates.}

Each product of a
matrix entry of \(g\) and a coefficient is a standard-tensor
coefficient. A one-coordinate integral shift can reach a compact
wall, at which its coherent character vanishes by compact Weyl
antisymmetry. It cannot cross the wall in that step. After \(j\) one-coordinate shifts an original opposite-sign
gap has decreased by at most \(j\). If its initial size is at least
\(h\), then throughout the expansion
\[
          \text{opposite-sign gap}\ge h-j>0\qquad(0\le j\le m<h).
\]
Thus the shifts cannot reach an opposite-sign wall. On the regular
summands the split Casimir has simple spectrum, by the difference
of the tensor and factor Casimir eigenvalues. The corresponding
projections are polynomials in first-order Lie operators of degree
at most \(n-1\).

The tensor form bounds the squared unitary norm
of a projected vector by its ordinary tensor Hilbert norm squared.
Commuting the projection with derived group actions shows that
this costs at most \(n-1\) derivatives in each seminorm. Induction
on \(m\), starting with unitarity, bounds a coefficient times any
monomial of \(m\) matrix entries by the two seminorms of order
\(m(n-1)\). Choose a matrix entry of maximal modulus. Since
\(\|g\|_{\rm HS}\le n\max_{a,b}|g_{ab}|\), the monomial
with \(m\) copies of this entry gives
\[
     |\langle v,\sigma(g)w\rangle|\ll_{m,n}
             \|g\|_{\rm HS}^{-m}\,p_m(v)p_m(w),
\]
where \(p_m\) uses derivatives through order \(m(n-1)\). Further fixed derivative orders change the
coefficient vectors and preserve the exponent. The constants may
depend on the coefficient and the derivatives.

\proofstep{Integrability on the period subgroup.}

Ambient decay yields convergence on the period subgroup only after integration in its unipotent coordinates.


The gap hypotheses give room for the exponent
\(a=h/(8\sqrt{2(r+1)})-1\), with \(m=\lceil a\rceil\). The
polynomial growth required for the period follows directly from
exponential coordinates on \(N\). Its dimension is
\(r(2k+2r+1)\le 2rn\), and the logarithm of a unipotent matrix is
polynomial of degree at most \(n-1\) in the entries minus the
identity.

In exponential coordinates that ball is contained in a ball
of radius \(O(M^{n-1})\). Hence its Haar volume is
\[
     \operatorname{vol}\{n\in N:\|n\|\le M\}
           \ll M^{(n-1)\dim N}\le M^{2rn^2}\qquad(M\ge1).
\]
The compact factor changes only the constant.
The numerical exponent in the period application exceeds this
power. The derivative constants suffice for the Fubini and cutoff
arguments because the output norm has the exact measure factor
\(c_{\rm meas}\). In the shift family Harish-Chandra's formal-degree
formula gives \(d_\nu=c_H\operatorname{Vand}(\lambda+\nu)/
\operatorname{Vand}(\rho_G)\) at fixed Haar. In particular, when
both periods occur, the normalization \(|s_\nu|=d_\nu\) gives
squared big-cell norm \(c_{\rm meas}d_\nu\).
\end{proof}
